\documentclass[11pt]{article}
\usepackage[a4paper,margin=2cm]{geometry}
\usepackage{amsmath,amssymb}
\usepackage[hidelinks]{hyperref}
\allowdisplaybreaks
\title{Qpu.General.Meaning.Prime}
\author{Tsvetan Rouschev}
\date{@uuidna/qpu v1.1.1, CC-BY-NC-ND-4.0}
\begin{document}
\maketitle
\noindent\url{https://github.com/uuidna/qpu/blob/v1.1.1/lean/Qpu/General/Meaning/Prime.lean}

\paragraph{\texttt{isPrime\_eq\_decide\_prime}}
`Nat.Prime p` means that `p` is a prime number, that is, a natural number at least 2 whose only divisors are `p` and `1`. The theorem `Nat.prime\_def` witnesses this description of a prime number.
\[ \forall (n : \mathbb{N}), \mathrm{Config.isPrime}\ n = \mathrm{decide} (\mathrm{Nat.Prime}\ n) \]
\noindent axioms: \texttt{propext}, \texttt{Classical.choice}, \texttt{Quot.sound}; proof: \url{https://github.com/uuidna/qpu/blob/v1.1.1/lean/Qpu/General/Meaning/Prime.lean\#L7}
\end{document}
